\chapter{The Higgs Scale, Running Couplings, and the Hierarchy Problem}

\lectureoverview{The final lecture begins by looking beyond the Standard
Model toward unification, supersymmetry, and the gauge-hierarchy problem.
To see why the Higgs field sits at the center of those questions, we first
reconstruct the Dirac Lagrangian and show that a fermion mass mixes left and
right chirality. The chiral weak interaction forbids that mass as a bare
term; a gauge-invariant Yukawa coupling supplies it only after the Higgs
field acquires a vacuum value. We then derive that vacuum value from the
Higgs potential, ask why every known elementary mass is tied to this one
small scale, and use superconductivity as a physical analog. After the
break, the lecture develops renormalization from conductor and dielectric
screening, carries the argument into QED and QCD, follows the gauge
couplings toward unification, and derives the Planck scale from
short-distance gravity. The electroweak scale then returns as the one
unexplained small number.}

\section{The Question Waiting Beyond the Standard Model}

The Standard Model is not presented here as a closed catalogue. At the
start of the evening I want to say where the next course will go. The
theory has difficulties rather than logical paradoxes: the gauge hierarchy,
the possible role of supersymmetry, and the question of unification. The
three gauge factors
\begin{equation}
  SU(3)_c\times SU(2)_L\times U(1)_Y
  \label{eq:l10-standard-model-group}
\end{equation}
describe strong, weak, and hypercharge interactions separately. Could they
be subgroups of one larger group? An often studied possibility is
\(SU(5)\), in which quarks and leptons occupy common multiplets. Whether
nature uses that particular construction is a different question. The
important clue is that the apparently separate structures fit remarkably
neatly into a larger mathematical one.

The Large Hadron Collider was built in part to expose whatever physics
settles these questions. In 2010 its most immediate quarry was the Higgs
boson. Finding a Higgs could never mean only adding one particle to a
table. A Higgs is identified by an interlocking set of properties: its
vacuum value gives masses to the weak vector bosons and fermions, while its
fluctuation couples to those same particles with definite strengths. A
resonance with the wrong couplings would not be the Standard Model Higgs.
\editorialnote{The lecture predates the 2012 discovery. The observed
approximately \(125\,\mathrm{GeV}\) scalar has, within present precision,
the spin, parity, and coupling pattern expected of the Standard Model
Higgs. No low-energy supersymmetric particles have been observed as of
2026. These later facts update the historical setting without changing the
argument developed in the room.}

The route into that argument begins not with the Higgs potential but with
the Dirac equation. We need to understand exactly what a fermion mass does
before asking where it comes from.

\section{From the Dirac Equation to Its Lagrangian}

\subsection{The older Dirac notation}

In the Hamiltonian notation used when the Dirac equation was first
introduced, its structure is
\begin{equation}
  i\partial_t\psi
  =
  \left(-i\alpha^i\partial_i+\beta m\right)\psi .
  \label{eq:l10-dirac-hamiltonian}
\end{equation}
Signs and factors of \(i\) depend on conventions, and the live board
temporarily suppresses some of them. The invariant content is the
separation between terms containing derivatives and a mass term containing
none.

\begin{figure}[htbp]
\centering
\includegraphics[width=0.82\textwidth]{lecture_10_figure_02.png}
\caption{The Dirac equation and Lagrangian in the older
\(\alpha,\beta\) notation. The board is being used structurally rather than
as a fixed sign convention.}
\lectureframe{00:05:52}
\label{fig:l10-dirac-alpha-beta}
\end{figure}

To obtain an action whose variation gives the Dirac equation, multiply the
equation by \(\psi^\dagger\). Moving every term to one side removes the need
for an equal sign, prompting the evening's small mathematical joke: once an
equation has been written as an expression that vanishes, the equal sign
has done its work. The corresponding first-order Lagrangian may be written
schematically as
\begin{equation}
  \mathcal L
  =
  \psi^\dagger
  \left(i\partial_t+i\alpha^i\partial_i-\beta m\right)\psi ,
  \label{eq:l10-dirac-old-lagrangian}
\end{equation}
with the understanding that a consistent metric convention fixes the
relative signs.

The derivative terms move an excitation from one spacetime point to a
nearby one and rotate its spinor components in the process. In a lattice
picture they are hopping terms. By contrast, the mass term absorbs and
re-emits the fermion at the same point. It has no derivative. This local
operation will shortly reveal itself as a repeated conversion between two
kinds of handedness.

\begin{classroomqa}
\audiencequestion{Should a \(\gamma^0\) be placed after
\(\psi^\dagger\)?\lecturetimestamp{00:05:59}}
\lecturerresponse{Not in the noncovariant form just written. But the
question points directly to the covariant notation. Define the Dirac
adjoint by \(\bar\psi=\psi^\dagger\gamma^0\); then the \(\gamma^0\) is
indeed built into the object that stands on the left.}
\end{classroomqa}

\subsection{The covariant form}

Set
\begin{equation}
  \bar\psi\equiv\psi^\dagger\beta
  =\psi^\dagger\gamma^0,
  \qquad
  \gamma^0\equiv\beta,
  \qquad
  \gamma^i\equiv\beta\alpha^i .
  \label{eq:l10-dirac-adjoint}
\end{equation}
Because \(\beta^2=1\), multiplication by \(\beta\) converts freely between
the two notations. The time derivative becomes
\(\bar\psi\gamma^0\partial_0\psi\), the spatial derivatives become
\(\bar\psi\gamma^i\partial_i\psi\), and the result is
\begin{equation}
  \boxed{\mathcal L_D
  =\bar\psi\left(i\gamma^\mu\partial_\mu-m\right)\psi .}
  \label{eq:l10-covariant-dirac}
\end{equation}

\begin{figure}[htbp]
\centering
\includegraphics[width=0.82\textwidth]{lecture_10_figure_03.png}
\caption{The covariant rewrite: \(\bar\psi=\psi^\dagger\gamma^0\),
\(\gamma^i=\beta\alpha^i\), and the compact Dirac Lagrangian.}
\lectureframe{00:10:17}
\label{fig:l10-dirac-covariant}
\end{figure}

This notation exposes the Lorentz character of each bilinear.
\(\bar\psi\psi\) is a scalar, whereas
\(\bar\psi\gamma^\mu\psi\) is a four-vector. The compact expression is not
merely prettier than the \(\alpha,\beta\) form; it tells us how the pieces
transform.

\begin{classroomqa}
\audiencequestion{Was the spinor at the far left of the compact expression
meant to be \(\bar\psi\)?\lecturetimestamp{00:09:39}}
\lecturerresponse{Yes. The missing bar is restored on the board. It matters:
\(\bar\psi\) is not \(\psi^\dagger\), but
\(\psi^\dagger\gamma^0\).}
\end{classroomqa}

\section{Chirality and What a Dirac Mass Actually Does}

Besides \(\gamma^0,\gamma^1,\gamma^2,\gamma^3\), define
\begin{equation}
  \gamma^5
  =i\gamma^0\gamma^1\gamma^2\gamma^3,
  \qquad
  (\gamma^5)^2=1 .
  \label{eq:l10-gamma-five}
\end{equation}
The factor of \(i\) shown here is the standard convention for the mostly
minus Minkowski metric. The eigenvalues are \(+1\) and \(-1\), and the
corresponding projectors are
\begin{equation}
  P_L=\frac{1-\gamma^5}{2},
  \qquad
  P_R=\frac{1+\gamma^5}{2},
  \qquad
  \psi_L=P_L\psi,\quad \psi_R=P_R\psi .
  \label{eq:l10-chiral-projectors}
\end{equation}

\begin{classroomqa}
\audiencequestion{Why is the extra matrix called \(\gamma^5\) when there is
no \(\gamma^4\)? And does its eigenvalue describe
helicity?\lecturetimestamp{00:10:42}}
\lecturerresponse{The name is historical: it is the additional matrix
formed from the product of the four spacetime gamma matrices. Its
eigenvalue labels chirality, not generally helicity. For a massless
particle the two notions can be identified appropriately, which is why the
language is easy to blur, but chirality is the Lorentz-representation label
needed by the weak gauge theory.}
\end{classroomqa}

The kinetic term preserves chirality. One useful way to see this is that
\(\{\gamma^\mu,\gamma^5\}=0\), so it separates into
\begin{equation}
  \bar\psi i\gamma^\mu\partial_\mu\psi
  =
  \bar\psi_L i\gamma^\mu\partial_\mu\psi_L
  +
  \bar\psi_R i\gamma^\mu\partial_\mu\psi_R .
  \label{eq:l10-chiral-kinetic}
\end{equation}
The scalar bilinear behaves differently:
\begin{equation}
  \bar\psi\psi
  =
  \bar\psi_L\psi_R+\bar\psi_R\psi_L .
  \label{eq:l10-chiral-mass}
\end{equation}

\begin{figure}[htbp]
\centering
\includegraphics[width=0.80\textwidth]{lecture_10_figure_04.png}
\caption{The board decomposes the scalar bilinear into left-to-right and
right-to-left pieces.}
\lectureframe{00:15:38}
\label{fig:l10-chiral-decomposition}
\end{figure}

Thus a Dirac mass is a local chirality-flipping operation:
\begin{equation}
  \psi_L\longleftrightarrow\psi_R .
\end{equation}
A massless fermion may propagate with definite chirality. A massive one
cannot remain in one chiral sector, because the mass term continually
mixes the two. It may seem surprising that this conversion has anything to
do with inertia, but it is precisely the relativistic field-theory form of
fermion rest mass.

\section{Why the Weak Symmetry Forbids a Bare Fermion Mass}

The charged weak current is purely left-chiral. The lecture first names the
\(Z\), then corrects the statement to the \(W^\pm\): \(Z\) couplings are
chiral and unequal, but not simply absent for every right-chiral fermion.
The clean statement is
\begin{equation}
  \mathcal L_{\mathrm{cc}}
  =
  -\frac{g}{\sqrt2}
  \left(\bar u_L\gamma^\mu d_L
  +\bar\nu_L\gamma^\mu e_L\right)W^+_\mu
  +\mathrm{h.c.}
  \label{eq:l10-charged-current}
\end{equation}
The left-handed fermions transform as \(SU(2)_L\) doublets; the
right-handed charged fermions transform as singlets.

Now compare those transformation rules with
\eqref{eq:l10-chiral-mass}. A term such as
\(\bar e_L e_R\) absorbs an object carrying one set of electroweak quantum
numbers and emits one carrying another. By itself it is not gauge
invariant. The familiar Dirac mass is therefore illegal as a bare term in
the electroweak Lagrangian.

This is the right order of explanation. Fermion masses are not first
inserted and then cosmetically associated with the Higgs. Gauge symmetry
first removes the obvious mass term. The Higgs field is then forced to
carry the missing quantum numbers.

\subsection{The Higgs as the charge carrier}

Let the Higgs transform as an \(SU(2)_L\) doublet with hypercharge
\(+\tfrac12\):
\begin{equation}
  H=
  \begin{pmatrix}
    H^+\\ H^0
  \end{pmatrix},
  \qquad
  \widetilde H=i\sigma^2H^* .
  \label{eq:l10-higgs-doublet}
\end{equation}
The lecture keeps the group indices schematic and calls the field
\(\phi\). In that notation, a Yukawa term has the form
\begin{equation}
  \mathcal L_Y
  =
  -G_Y\left(
    \bar\psi_L\phi\psi_R+
    \bar\psi_R\phi^\dagger\psi_L
  \right).
  \label{eq:l10-yukawa-schematic}
\end{equation}
The Higgs field carries away the weak charge that a naked chirality flip
would fail to conserve.

The complete one-generation Standard Model expression makes the same
bookkeeping explicit:
\begin{equation}
  \mathcal L_Y
  =
  -y_d\bar Q_LHd_R
  -y_u\bar Q_L\widetilde H u_R
  -y_e\bar L_LHe_R
  +\mathrm{h.c.}
  \label{eq:l10-yukawa-full}
\end{equation}
Each term is gauge invariant, and each fermion species carries its own
dimensionless Yukawa coupling.

\begin{figure}[htbp]
\centering
\includegraphics[width=0.80\textwidth]{lecture_10_figure_05.png}
\caption{The Higgs-assisted chirality flip and the symmetry-breaking
potential appear together on the board.}
\lectureframe{00:18:54}
\label{fig:l10-yukawa-potential-board}
\end{figure}

\section{One Interaction Becomes Both Mass and Higgs Coupling}

Nature chooses a vacuum in which the neutral Higgs component does not
vanish. In the lecture's schematic real-field notation,
\begin{equation}
  \phi=F+\mathcal H ,
  \label{eq:l10-f-plus-h}
\end{equation}
where \(F\) is the nonfluctuating vacuum shift and \(\mathcal H\) is the
propagating Higgs fluctuation. Substitution into
\eqref{eq:l10-yukawa-schematic} gives
\begin{align}
  \mathcal L_Y
  ={}&
  -G_YF
  \left(\bar\psi_L\psi_R+\bar\psi_R\psi_L\right)
  \notag\\
  &-G_Y\mathcal H
  \left(\bar\psi_L\psi_R+\bar\psi_R\psi_L\right).
  \label{eq:l10-yukawa-split}
\end{align}
The first line is a Dirac mass and the second is an interaction with a
physical Higgs quantum:
\begin{equation}
  m_f=G_YF
  \qquad\hbox{in the lecture's normalization.}
  \label{eq:l10-mass-lecture-normalization}
\end{equation}

In the conventional Standard Model normalization,
\begin{equation}
  H(x)=\frac{1}{\sqrt2}
  \begin{pmatrix}
    0\\ v+h(x)
  \end{pmatrix},
  \qquad
  m_f=\frac{y_fv}{\sqrt2},
  \qquad
  \mathcal L_{hff}=-\frac{m_f}{v}h\bar f f,
  \label{eq:l10-standard-higgs-normalization}
\end{equation}
with \(v\simeq246\,\mathrm{GeV}\). The lecture's rough
\(F\sim200\,\mathrm{GeV}\) is this same electroweak scale with normalization
details suppressed.

The measured \(W\) and \(Z\) masses determine the scale \(F\). The measured
fermion masses then determine the Yukawa couplings. The electron's coupling
is tiny; the top quark's is of order unity. If this were all the mechanism
did, we would have put in one unknown number per fermion and read the same
information back out.

But the second line of \eqref{eq:l10-yukawa-split} is new experimental
content. It says that a fermion can emit or absorb a Higgs and, read in the
crossed channel, that a Higgs can decay into a fermion-antifermion pair.
Schematically,
\begin{equation}
  \mathcal M(h\to f\bar f)\propto y_f,
  \qquad
  \Gamma(h\to f\bar f)
  =
  N_c\frac{y_f^2m_h}{16\pi}
  \left(1-\frac{4m_f^2}{m_h^2}\right)^{3/2},
  \label{eq:l10-higgs-fermion-width}
\end{equation}
at tree level. Here \(N_c=3\) for quarks and \(1\) for leptons. A larger
fermion mass means a larger Yukawa coupling, provided the decay is
kinematically allowed. This does not mean every heavy channel dominates:
color factors, phase space, loop-induced channels, and decays to vector
bosons also matter. It means the fermionic coupling itself is proportional
to mass.

\begin{classroomqa}
\audiencequestion{How would one actually predict a Higgs decay time without
doing every technical detail?\lecturetimestamp{00:25:32}}
\lecturerresponse{The interaction supplies a Hamiltonian matrix element.
Draw the corresponding Feynman diagram, calculate its transition
amplitude, square the amplitude, and sum or integrate over all allowed final
states. The result is the decay rate. The lifetime is the inverse total
width, \(\tau=\hbar/\Gamma_{\mathrm{tot}}\). This is the same quantum
mechanical logic used in particle calculations since the early development
of quantum field theory.}
\end{classroomqa}

At the time of the lecture, this Higgs phenomenology was still prediction
rather than direct observation. The \(W\) and \(Z\) self-interactions had
already provided stringent tests of the gauge structure; the Higgs and its
direct couplings remained to be found. The historical uncertainty belongs
in the record. So does the later outcome: a Higgs-like particle was found,
and measurements of its couplings are now tests of the same mass-generating
structure developed here.

\section{The Potential That Chooses F}

The vacuum shift is not chosen by hand after writing the Yukawa coupling.
It is determined by the Higgs potential. For one real cross-section through
the full doublet, take
\begin{equation}
  V(\phi)
  =
  -\frac{\mu^2}{2}\phi^2
  +\frac{\lambda}{4}\phi^4,
  \qquad
  \mu^2>0,\quad\lambda>0.
  \label{eq:l10-real-higgs-potential}
\end{equation}
Near \(\phi=0\), the negative quadratic term makes an upside-down
parabola. Calling \(\mu\) an imaginary mass is shorthand for saying that
the origin is unstable; it does not mean that the physical Higgs
excitation about the true vacuum has an imaginary mass.

The quadratic term alone would drive the field without bound and leave no
lowest-energy state. The positive quartic term eventually dominates and
turns the potential upward. An odd power would spoil the toy
\(\phi\mapsto-\phi\) symmetry. A quartic is the first even stabilizing
interaction and, in four spacetime dimensions, the highest polynomial
scalar interaction with a dimensionless coupling.

Differentiate:
\begin{align}
  \frac{dV}{d\phi}
  &=-\mu^2\phi+\lambda\phi^3
  \notag\\
  &=\phi\left(-\mu^2+\lambda\phi^2\right).
  \label{eq:l10-potential-derivative}
\end{align}
The stationary point \(\phi=0\) is the unstable top. The two minima satisfy
\begin{equation}
  \boxed{F^2=\frac{\mu^2}{\lambda}},
  \qquad
  F=\frac{\mu}{\sqrt\lambda}
  \label{eq:l10-vacuum-value}
\end{equation}
for this real-field convention.

\begin{figure}[htbp]
\centering
\includegraphics[width=0.84\textwidth]{lecture_10_figure_06.png}
\caption{The completed minimization on the board:
\(V=-\mu^2\phi^2/2+\lambda\phi^4/4\),
\(V'=0\), and \(\phi^2=\mu^2/\lambda\).}
\lectureframe{00:34:32}
\label{fig:l10-higgs-minimization}
\end{figure}

The factors \(1/2\) and \(1/4\) were inserted precisely so they disappear
upon differentiation. Near the origin \(\phi^4\) is smaller than
\(\phi^2\); at large field it grows faster and stabilizes the system. That
simple competition produces the double well.

The dimensions also settle a momentary verbal slip in the transcript.
\(\lambda\) is dimensionless in four spacetime dimensions, while \(\mu\)
and \(F\) have dimensions of mass. Expanding \(\phi=F+\mathcal H\), the
curvature of the potential gives
\begin{equation}
  m_{\mathcal H}^2
  =
  \left.\frac{d^2V}{d\phi^2}\right|_{\phi=F}
  =
  -\mu^2+3\lambda F^2
  =
  2\mu^2
  =
  2\lambda F^2.
  \label{eq:l10-higgs-curvature}
\end{equation}
In the full Standard Model notation,
\(V(H)=-\mu^2H^\dagger H+\lambda(H^\dagger H)^2\), one likewise has
\(m_h^2=2\lambda v^2\). The measured Higgs mass and \(v\) imply
\(\lambda\) near \(0.13\) at the electroweak scale. The lecture's 2010
estimate was intentionally rough; the structural claim was that
\(\lambda\) is not an absurdly tiny number.

\section{Why Is Everything We See Tied to This Scale?}

The same \(F\) controls the \(W\) and \(Z\) masses and, multiplied by
different Yukawa couplings, all charged-fermion masses. That empirical
concentration raises a more interesting question than the algebra:
why should all known elementary particles borrow their masses from the
same symmetry-breaking phenomenon?

There is no theorem that every possible particle must. Imagine a
vectorlike fermion whose left- and right-chiral components transform in the
same way under the gauge group. Then
\(\bar\psi_L\psi_R+\mathrm{h.c.}\) can be gauge invariant without a Higgs.
A complete gauge singlet can likewise carry a mass unrelated to
electroweak breaking. Such particles could naturally be much heavier than
the particles tied to \(F\). The particles we have discovered may be light
enough to see precisely because symmetry would make them massless when
\(F\) vanishes.

The electroweak scale is therefore the exceptional number. The speculative
picture is that generic unprotected masses lie near a much larger
fundamental scale, while particles protected by chiral gauge symmetry are
dragged down with \(F\).

\begin{classroomqa}
\audiencequestion{Is the Planck mass one candidate for the fundamental
mass scale with which \(F\) should be compared?\lecturetimestamp{00:38:43}}
\lecturerresponse{Yes. Another is the scale suggested by gauge-coupling
unification. Both will emerge from arguments later in the lecture, even
though no particles of those enormous masses have been directly observed.}
\end{classroomqa}

\begin{classroomqa}
\audiencequestion{What does it mean to say that the other scales are about
fifteen orders of magnitude above \(F\), and what particles correspond to
them?\lecturetimestamp{00:39:54}}
\lecturerresponse{The electroweak scale is of order
\(10^2\,\mathrm{GeV}\), while unification is suggested near
\(10^{16}\,\mathrm{GeV}\) and the Planck scale lies near
\(10^{19}\,\mathrm{GeV}\). We have not produced particles at those masses.
We infer the scales indirectly: one from extrapolating measured gauge
couplings and the other from combining quantum mechanics, relativity, and
Newton's constant.}
\end{classroomqa}

\section{A Charged Condensate in the Laboratory}

Before the break, an audience question asks for an ordinary condensed
matter analog of the Higgs mechanism. The answer develops in stages rather
than by analogy alone.

\begin{classroomqa}
\audiencequestion{Beyond ultracold Bose condensates, is there a common
solid-state example of a Higgs-like mechanism?\lecturetimestamp{00:40:50}}
\lecturerresponse{Yes. A neutral bosonic field can acquire a nonzero
expectation value in a Bose condensate or superfluid. If the condensed
object is charged, the same basic phenomenon alters the gauge field.
Superconductivity supplies the familiar example: charged Cooper pairs
condense, and electromagnetic fields acquire a finite penetration length
inside the material.}
\end{classroomqa}

At molecular resolution a helium atom can be treated as one effective
bosonic particle with a field \(\Phi_{\mathrm{He}}\). Empty space
corresponds to zero field; a superfluid phase has
\(\langle\Phi_{\mathrm{He}}\rangle\ne0\) and spontaneously breaks the
relevant global phase symmetry. Helium atoms are neutral, so this is not
yet the gauge-field mass mechanism.

In a superconductor, electrons form weakly bound Cooper pairs. The pair
field is schematically an electron bilinear,
\begin{equation}
  \Delta(x)\sim\langle\psi_\downarrow(x)\psi_\uparrow(x)\rangle,
  \label{eq:l10-cooper-order-parameter}
\end{equation}
and carries electric charge \(2e\). When \(\Delta\) condenses, the
electromagnetic gauge field responds exactly as a gauge field coupled to a
charged order parameter must. Static magnetic fields decay over the London
penetration depth:
\begin{equation}
  B(x)\propto e^{-x/\lambda_L}.
  \label{eq:l10-meissner-screening}
\end{equation}
The inverse penetration length plays the role of an effective photon mass
inside the medium. It is not a vacuum photon mass; it is a property of the
superconducting phase.

\begin{classroomqa}
\audiencequestion{Are the two electrons in a Cooper pair not separated,
even lying on opposite sides of the Fermi
surface?\lecturetimestamp{00:43:13}}
\lecturerresponse{They are not a tiny molecule localized at one point. A
Cooper pair is extended and is naturally described in momentum space by
correlated states near opposite momenta. For the symmetry argument,
however, the pair acts as a single bosonic charged order parameter. Its
condensation is what matters.}
\end{classroomqa}

Historically this relationship is more than a teaching metaphor.
Superconductivity and Anderson's analysis of gauge fields in a broken
phase helped make the relativistic Higgs mechanism intelligible. The
particle-physics construction was developed independently in relativistic
field theory, but the shared mathematics is real.

\section{Bare Parameters and Measured Parameters}

After the break, we return to the Lagrangian. It contains
parameters called masses and charges, but those input parameters are rarely
identical to quantities measured without qualification. High-frequency
modes interact with low-frequency modes and dress them. The systematic
relation between parameters defined at one resolution and observables at
another is renormalization.

Electric charge gives a picture that can be built without calculation.
At large separation, Coulomb's law defines the observed charge:
\begin{equation}
  V(r)\sim\frac{e_1e_2}{4\pi r},
  \qquad
  F(r)\sim\frac{e_1e_2}{4\pi r^2}.
  \label{eq:l10-coulomb-law}
\end{equation}
The lecture suppresses \(4\pi\) and writes \(e^2/r\) or \(e^2/r^2\).
What matters is that the coefficient inferred from a force measurement can
depend on the distance at which the system is probed.

\subsection{A conductor: complete screening}

Place a positive test charge in an ideal conductor. Its electric field
drives mobile negative charge toward it and positive charge away. Current
flows until the local source is surrounded by an opposite cloud and no
field remains in the bulk outside that cloud; compensating charge is pushed
to the remote boundary. A negative source produces the opposite
rearrangement. Two dressed sources far apart can then exert essentially no
force, even though the charges inserted microscopically were nonzero.

\begin{figure}[htbp]
\centering
\includegraphics[width=0.82\textwidth]{lecture_10_figure_07.png}
\caption{A source charge surrounded by an oppositely charged screening
cloud in the conductor discussion.}
\lectureframe{00:50:24}
\label{fig:l10-conductor-screening}
\end{figure}

\begin{classroomqa}
\audiencequestion{Is this rearrangement what is meant by
screening?\lecturetimestamp{00:51:25}}
\lecturerresponse{Yes. At distances larger than the cloud, the inserted
charge is screened. Bring another probe inside the cloud and some of the
screening can no longer lie between source and probe; the underlying charge
then appears stronger.}
\end{classroomqa}

The cloud has a characteristic scale set by the dynamics and density of
the mobile charges. At separations \(r\) smaller than that scale, the probe
penetrates the cloud. One may therefore define a running charge \(e(r)\)
through
\begin{equation}
  F(r)\equiv\frac{e^2(r)}{4\pi r^2}.
  \label{eq:l10-running-charge-definition}
\end{equation}
For the idealized conductor, \(e(r)\) approaches zero at very long
distance and the microscopic value at very short distance.

\subsection{A dielectric: partial screening}

In an insulator, electrons are bound to atoms rather than free to cross the
sample. A crude but useful model pictures each electron cloud attached to
its ion by a spring. Without an external field, the average positive and
negative charge centers coincide. An applied field shifts them slightly
apart and induces a dipole.

Around a negative source, nearby electron clouds shift away and leave
their positive ions slightly closer; around a positive source they shift
the other way. The bound charges cannot flow far enough to neutralize the
source completely, so the long-distance field is reduced by the dielectric
constant rather than eliminated:
\begin{equation}
  V_{\mathrm{medium}}(r)
  =
  \frac{1}{\epsilon_r}
  \frac{e_1e_2}{4\pi\epsilon_0r}.
  \label{eq:l10-dielectric-coulomb}
\end{equation}
At distances inside the polarization cloud, the probe again sees more of
the bare source. Complete and partial screening are two versions of the
same scale-dependent idea.

\section{The QED Vacuum as a Polarizable Medium}

Quantum electrodynamics has no material atoms in its vacuum, but it has
charged quantum fields. Electron-positron fluctuations contribute loops to
photon propagation. In the older Dirac-sea language, the electric field
distorts the filled negative-energy states; in modern language, vacuum
polarization dresses the photon propagator. The medium picture says that a
virtual pair is polarized: its positive and negative components shift in
opposite directions.

A positive source is surrounded by a small excess of negative
polarization charge, and a negative source by positive polarization charge.
The charge measured far away is therefore smaller than the parameter one
would associate with an unresolved source. Unlike an ordinary material,
the vacuum contains fluctuations over a continuum of momenta. Low-energy
pairs correspond to longer distance structure; high-energy pairs probe
shorter distances. There is no final cloud radius at which the process
simply stops.

As two electrons are brought closer together, the largest polarization
clouds cease to fit between them first. Bringing them closer peels away
successively smaller clouds. The effective electric coupling grows
logarithmically with inverse distance, or equivalently with energy:
\begin{equation}
  \alpha(Q)
  \simeq
  \frac{\alpha(\mu)}
  {1-\dfrac{2\alpha(\mu)}{3\pi}
  \ln\!\left(\dfrac{Q}{\mu}\right)}
  \label{eq:l10-qed-running}
\end{equation}
for one charged Dirac fermion well above its threshold at leading order.
Additional charged species change the coefficient when their thresholds
are crossed.

\begin{figure}[htbp]
\centering
\includegraphics[width=0.82\textwidth]{lecture_10_figure_08.png}
\caption{The QED effective charge rises slowly as the horizontal scale
\(1/R\), and therefore the probing energy, increases.}
\lectureframe{01:02:13}
\label{fig:l10-qed-running-board}
\end{figure}

Electron scattering gives the operational version. Higher collision
energy means shorter de Broglie wavelength and smaller closest approach.
The amplitude must therefore use the coupling appropriate to that momentum
transfer, not one fixed number for every scale.

\begin{classroomqa}
\audiencequestion{At what distance does the running first become
measurable?\lecturetimestamp{01:03:54}}
\lecturerresponse{Electron vacuum polarization becomes relevant around the
electron Compton scale, \(\lambda_C=\hbar/(m_ec)\), roughly
\(4\times10^{-11}\,\mathrm{cm}\) for the reduced wavelength. The effect is
small and logarithmic, so precision is required. Heavier charged particles
join the running at their own thresholds.}
\end{classroomqa}

\begin{classroomqa}
\audiencequestion{Does speaking of the smallest distance assume that we
already know what happens there?\lecturetimestamp{01:04:49}}
\lecturerresponse{No. We extrapolate the theory only while its particle
content and interactions remain valid. New physics can change the running.
Experiment does not tell us that QED remains the complete description to
arbitrarily short distance.}
\end{classroomqa}

\begin{classroomqa}
\audiencequestion{For gluons, is the scale dependence the opposite? Why?
\lecturetimestamp{01:05:03}}
\lecturerresponse{Yes. Quark loops screen color charge in a way analogous
to charged matter, but gluons themselves carry color and interact
nonlinearly. Their contribution has the opposite sign and dominates for
the observed number of quark flavors. The result is antiscreening and
asymptotic freedom.}
\end{classroomqa}

\begin{classroomqa}
\audiencequestion{How could one determine a bare electric charge if doing
so requires an infinitely small distance?\lecturetimestamp{01:05:28}}
\lecturerresponse{One does not measure a unique regulator-independent bare
charge directly. One defines a parameter at a chosen short-distance cutoff
or renormalization scale, then adjusts it so a measured low-energy
quantity remains fixed. The scale-dependent renormalized coupling is the
useful physical object.}
\end{classroomqa}

\begin{classroomqa}
\audiencequestion{Does the QED charge approach a finite value at infinite
energy?\lecturetimestamp{01:05:39}}
\lecturerresponse{Not in perturbative QED extrapolated by itself. Its
positive beta function makes the coupling continue to grow logarithmically
and leads formally to a Landau pole at a fantastically high scale. Long
before treating that formal pole as physical, one expects the description
to be embedded in electroweak theory and ultimately altered by further
physics.}
\end{classroomqa}

\subsection{What renormalization changes and what it preserves}

Choose a cutoff length \(a\), call the corresponding input \(e(a)\), and
calculate a fixed long-distance observable. A different cutoff \(a'\) is
equally acceptable if the input is changed to \(e(a')\) so the same
observable is reproduced:
\begin{equation}
  \frac{d}{d\ln\mu}\,\mathcal O_{\mathrm{physical}}=0,
  \qquad
  \mu\frac{de}{d\mu}=\beta(e).
  \label{eq:l10-rg-invariance}
\end{equation}
That compensating flow is the renormalization group.

\begin{classroomqa}
\audiencequestion{Must the word bare refer specifically to the Planck
length?\lecturetimestamp{01:06:20}}
\lecturerresponse{No. One may terminate the description at any chosen
cutoff and call the parameter there bare. Changing the cutoff while
changing that parameter appropriately leaves accessible physics
unchanged. The Planck length is a plausible place where ordinary quantum
field theory must be replaced, not a required convention for
renormalization.}
\end{classroomqa}

\begin{classroomqa}
\audiencequestion{Is this procedure a matter of zooming in and attenuating
the description?\lecturetimestamp{01:07:15}}
\lecturerresponse{It is zooming in conceptually. We ask how the parameters
of a finer-grained description must change so that predictions at a fixed,
experimentally accessible wavelength remain the same.}
\end{classroomqa}

\begin{classroomqa}
\audiencequestion{Does the same running phenomenon happen with
gravity?\lecturetimestamp{01:08:14}}
\lecturerresponse{Gravity has scale dependence too, but not simply the same
dielectric-screening mechanism. For the present lecture a dimensional
argument is enough: localization raises the energy, energy gravitates, and
the resulting dimensionless gravitational strength grows rapidly with
energy.}
\end{classroomqa}

\section{QCD: Antiscreening, Flux Tubes, and Asymptotic Freedom}

At hadronic distances, chromoelectric field lines do not spread
spherically like an electromagnetic Coulomb field. Gluon self-interactions
squeeze them into a narrow tube between a quark and antiquark. If the tube
has approximately fixed energy per unit length \(\sigma\), then
\begin{equation}
  V_{q\bar q}(r)
  \simeq
  -\frac{C_F\alpha_s(r)}{r}
  +\sigma r .
  \label{eq:l10-cornell-potential}
\end{equation}
At large \(r\), the linear term dominates. Stretching the pair creates more
of the tube, like pulling a strand of taffy, and the energy grows rather
than tending to zero. This is the confinement side of the story.

\begin{figure}[htbp]
\centering
\includegraphics[width=0.84\textwidth]{lecture_10_figure_09.png}
\caption{Chromoelectric field lines are drawn as a tube joining separated
color sources rather than as a freely spreading Coulomb field.}
\lectureframe{01:09:56}
\label{fig:l10-qcd-flux-tube-board}
\end{figure}

At distances shorter than the tube radius, the potential returns
approximately to a Coulomb form, but its coupling decreases with energy.
At one loop,
\begin{equation}
  \alpha_s(Q)
  =
  \frac{4\pi}
  {\left(11-\frac23n_f\right)
  \ln(Q^2/\Lambda_{\mathrm{QCD}}^2)} .
  \label{eq:l10-asymptotic-freedom}
\end{equation}
For \(n_f<17\), the coefficient is positive and
\(\alpha_s(Q)\to0\) as \(Q\to\infty\). QCD is weak at very short distance
and strong at long distance. This is the precise running-coupling version
of the lecture's opposite behavior.

\begin{classroomqa}
\audiencequestion{What happens to the weak coupling between the QED and QCD
curves?\lecturetimestamp{01:11:53}}
\lecturerresponse{Its running is intermediate and depends on which
electroweak factor is being followed. Above the electroweak scale the
proper quantities are the \(U(1)_Y\) and \(SU(2)_L\) couplings. One rises
and the other falls slowly. Together with the falling strong coupling,
their extrapolated values come surprisingly close at high energy.}
\end{classroomqa}

\section{Running Toward Unification}

Measure the three gauge couplings near the weak scale and evolve them with
the particle content of a proposed theory. Because every beta function is
logarithmic, an enormous change in energy produces only a moderate change
in coupling. The Standard Model trajectories approach one another but do
not meet exactly. With the low-energy particle content of the minimal
supersymmetric Standard Model, the one-loop trajectories meet much more
closely near
\begin{equation}
  M_{\mathrm{GUT}}\sim10^{16}\,\mathrm{GeV}.
  \label{eq:l10-gut-scale}
\end{equation}

The plot is evidence, not proof. Thresholds from undiscovered particles,
higher-loop corrections, and details of a unified theory shift the curves.
The 2010 lecture emphasizes the approximately percent-level convergence in
a supersymmetric spectrum because it made weak-scale supersymmetry
especially attractive. The absence of superpartners at collider energies
now means that the simplest version of that spectrum is not realized as
originally expected; precision unification remains model dependent.

\begin{figure}[htbp]
\centering
\includegraphics[width=0.84\textwidth]{lecture_10_figure_10.png}
\caption{The board's qualitative coupling plot: the strong curve falls,
the electroweak curves evolve slowly toward it, and gravity rises steeply
near the far ultraviolet.}
\lectureframe{01:16:29}
\label{fig:l10-unification-board}
\end{figure}

Why regard the vicinity of a meeting point as fundamental? It is where the
input may become simplest. Three unrelated low-energy couplings could be
different dressed descendants of one high-energy coupling associated with
a single gauge group. This is the numerical clue behind the earlier
\(SU(5)\) question.

Gravity appears on the same logarithmic sketch as a curve that rises much
faster. It does not cross the gauge curves at exactly the GUT scale, but it
becomes order one within a few decades of energy. On a logarithmic axis,
that proximity is suggestive: \(10^{16}\), \(10^{17}\), \(10^{18}\), and
\(10^{19}\,\mathrm{GeV}\) are vastly closer to one another than any is to
the electroweak scale.

\section{Why Gravity Strengthens at Short Distance}

\subsection{Energy, not rest mass alone, gravitates}

Newton's force between two slowly moving masses is
\begin{equation}
  F_G=\frac{Gm_1m_2}{r^2}.
  \label{eq:l10-newton-force}
\end{equation}
For a relativistic dimensional estimate, replace each rest mass by the
corresponding energy divided by \(c^2\):
\begin{equation}
  F_G(r)\sim\frac{GE_1E_2}{c^4r^2}.
  \label{eq:l10-relativistic-gravity-force}
\end{equation}
This is not a substitute for general relativity; stress, momentum flow, and
pressure also source gravity. It is the correct scaling estimate for the
question at hand.

Localize each particle to a region of size \(r\). The uncertainty
principle requires
\begin{equation}
  p\sim\frac{\hbar}{r}.
  \label{eq:l10-localization-momentum}
\end{equation}
At sufficiently small \(r\), the momentum is ultrarelativistic, so
\begin{equation}
  E\simeq pc\sim\frac{\hbar c}{r}.
  \label{eq:l10-localization-energy}
\end{equation}
Substituting two such energies into
\eqref{eq:l10-relativistic-gravity-force} gives
\begin{equation}
  F_G(r)
  \sim
  \frac{G\hbar^2}{c^2r^4}.
  \label{eq:l10-short-distance-gravity-force}
\end{equation}
The corresponding potential-energy scale is
\begin{equation}
  U_G(r)\sim-\frac{G\hbar^2}{c^2r^3}.
  \label{eq:l10-short-distance-gravity-potential}
\end{equation}
This is the distinction corrected during the live calculation:
the potential scales as \(1/r^3\), while the force scales as \(1/r^4\).

\begin{classroomqa}
\audiencequestion{For an ultrarelativistic particle, is the energy
\(p^2/(2m)\)?\lecturetimestamp{01:22:18}}
\lecturerresponse{No. That is the nonrelativistic formula. From
\(E^2=p^2c^2+m^2c^4\), the limit \(pc\gg mc^2\) gives
\(E\simeq pc\). Combined with \(p\sim\hbar/r\), the localization energy is
\(\hbar c/r\).}
\end{classroomqa}

\begin{figure}[htbp]
\centering
\includegraphics[width=0.84\textwidth]{lecture_10_figure_11.png}
\caption{The board substitutes relativistic energies into Newton's law
while checking the powers of \(c\) and \(r\).}
\lectureframe{01:23:29}
\label{fig:l10-gravity-derivation-board}
\end{figure}

\begin{classroomqa}
\audiencequestion{Did the factors of \(c\) cancel, and was the claimed
\(1/r^3\) law a force law?\lecturetimestamp{01:23:38}}
\lecturerresponse{The board calculation is corrected in real time. The two
effective masses contribute \(E_1E_2/c^4\); with
\(E_i\sim\hbar c/r\), two powers of \(c\) remain in the denominator. The
resulting force is \(G\hbar^2/(c^2r^4)\). The \(1/r^3\) expression belongs
to the potential energy.}
\end{classroomqa}

\begin{classroomqa}
\audiencequestion{The gravitational force under discussion is attractive,
correct?\lecturetimestamp{01:24:42}}
\lecturerresponse{Yes. The equations above display magnitudes. The
Newtonian gravitational potential carries a minus sign for positive
energies, so the force is attractive in this estimate.}
\end{classroomqa}

\begin{classroomqa}
\audiencequestion{Is gravity gaining relative strength because
localization raises energy, while electric charge itself does not gain the
same kinematic factor?\lecturetimestamp{01:25:02}}
\lecturerresponse{Exactly. Energy is the gravitational charge. Squeezing a
quantum state raises its momentum and energy, so it also raises what
sources gravity. Electric charge is dimensionless and does not acquire
that direct \(1/r\) factor, apart from its much slower logarithmic
renormalization.}
\end{classroomqa}

\subsection{The Planck crossing}

The electromagnetic force between charges of order one in natural units
has magnitude
\begin{equation}
  F_{\mathrm{EM}}(r)
  \sim
  \alpha\,\frac{\hbar c}{r^2}.
  \label{eq:l10-em-force-alpha}
\end{equation}
Equating it to \eqref{eq:l10-short-distance-gravity-force} gives
\begin{equation}
  r^2\sim\frac{G\hbar}{\alpha c^3}.
  \label{eq:l10-gravity-em-crossing}
\end{equation}
The lecture sets the dimensionless gauge factor to order one to expose the
fundamental gravitational length:
\begin{equation}
  \boxed{\ell_P=\sqrt{\frac{\hbar G}{c^3}}}
  \simeq1.616\times10^{-35}\,\mathrm{m}.
  \label{eq:l10-planck-length}
\end{equation}
Keeping the actual \(\alpha\) moves the equality with electromagnetism by
an order-one logarithmic-scale factor, \(r\sim\ell_P/\sqrt\alpha\). The
fundamental statement is cleaner in natural units:
\begin{equation}
  \alpha_G(E)\equiv\frac{GE^2}{\hbar c^5}
  =\left(\frac{E}{E_P}\right)^2,
  \qquad
  E_P=\sqrt{\frac{\hbar c^5}{G}}
  \simeq1.22\times10^{19}\,\mathrm{GeV}.
  \label{eq:l10-gravitational-coupling}
\end{equation}
Gravity becomes strong when \(E\) approaches \(E_P\).

\begin{classroomqa}
\audiencequestion{Was \(G\) or the electric coupling being set equal to one
in the comparison?\lecturetimestamp{01:27:02}}
\lecturerresponse{The electric coupling was being treated as an order-one
number for the rough estimate. Newton's constant must remain: it is the
dimensional quantity that determines the Planck scale.}
\end{classroomqa}

\begin{classroomqa}
\audiencequestion{The intermediate powers of \(r\), \(\hbar\), and \(c\)
do not appear dimensionally consistent. What is the corrected
calculation?\lecturetimestamp{01:27:47}}
\lecturerresponse{Use energies from the start. Localization gives
\(E\sim\hbar c/r\), hence
\(U_G\sim GE^2/(c^4r)=G\hbar^2/(c^2r^3)\). Compare this with
\(U_{\mathrm{EM}}\sim\alpha\hbar c/r\). Setting
\(\alpha\) to order one yields
\(r^2=G\hbar/c^3\). The class and board arrive at
\(r=\sqrt{G\hbar/c^3}\), the Planck length.}
\end{classroomqa}

\begin{figure}[htbp]
\centering
\includegraphics[width=0.84\textwidth]{lecture_10_figure_12.png}
\caption{The corrected endpoint of the live dimensional check:
\(r=\sqrt{\hbar G/c^3}\).}
\lectureframe{01:32:16}
\label{fig:l10-planck-length-board}
\end{figure}

\section{The Hierarchy Problem Returns}

The gauge couplings become comparable near a very high scale, and gravity
becomes strong near the Planck scale. These enormous scales now look less
like arbitrary large numbers and more like places where a microscopic
description might simplify. Against them, the electroweak vacuum value
\(F\) is anomalously small:
\begin{equation}
  \frac{v}{E_P}\sim2\times10^{-17}
  \label{eq:l10-electroweak-planck-ratio}
\end{equation}
in units with \(c=1\). Using a GUT scale instead of the Planck scale
changes the precise ratio but not the puzzle.

Once \(F\) is small, the masses of the \(W\), \(Z\), charged leptons, and
quarks follow because gauge symmetry makes their masses proportional to
it. Particles without that protection can remain heavy. The outstanding
question is therefore not why every known mass is independently tiny.
It is why the Higgs sector supplies this one tiny symmetry-breaking scale.

\begin{classroomqa}
\audiencequestion{If one could calculate \(F\), would all the other masses
then be calculable?\lecturetimestamp{01:34:39}}
\lecturerresponse{It would determine the overall electroweak scale and,
with the gauge and Yukawa couplings supplied, the corresponding particle
masses. The remaining problem is that an ordinary ultraviolet-sensitive
calculation tends to drive the Higgs mass parameter toward the high
fundamental scale unless a symmetry or dynamical mechanism protects it.}
\end{classroomqa}

The word fine-tuned anticipates that sensitivity. In a cutoff description,
a scalar mass parameter receives corrections of the rough form
\begin{equation}
  \delta\mu^2
  \sim
  \frac{C}{16\pi^2}\Lambda^2,
  \label{eq:l10-higgs-quadratic-sensitivity}
\end{equation}
where \(C\) combines gauge, Yukawa, and scalar couplings and \(\Lambda\)
marks the scale at which the effective theory is replaced. If
\(\Lambda\) is near a GUT or Planck scale, obtaining
\(\mu^2\sim v^2\) appears to require a very precise cancellation between
the input and radiative contributions. This is the naturalness form of the
gauge-hierarchy problem. It is an editorial completion of the calculation
the lecture promises for the following quarter, not a claim that a
particular regulator artifact is itself observable.

\begin{classroomqa}
\audiencequestion{Should the Planck length be visualized as the final
screening-cloud radius?\lecturetimestamp{01:35:03}}
\lecturerresponse{No. The dielectric picture explains gauge-charge
renormalization, but the Planck argument used localization energy and the
dimensional gravitational coupling. There is no literal gravitational
screening cloud in this derivation.}
\end{classroomqa}

\subsection{The closing questions about interactions}

The last minutes return to the Higgs potential and to what terms a
Lagrangian may contain.

\begin{classroomqa}
\audiencequestion{Could the Higgs potential include
\(\phi^6\) and still remain symmetric?\lecturetimestamp{01:35:20}}
\lecturerresponse{Yes. Even powers preserve the toy
\(\phi\mapsto-\phi\) symmetry, whereas \(\phi^5\) would break it. In four
dimensions the coefficient of \(\phi^6\) has negative mass dimension, so
the term is suppressed by a high scale, for example
\(c_6\phi^6/\Lambda^2\). It is a valid effective interaction even though it
is not power-counting renormalizable.}
\end{classroomqa}

\begin{classroomqa}
\audiencequestion{Do the powers of the Higgs field represent creation and
annihilation processes, so that \(\phi^4\) predicts Higgs-Higgs
scattering?\lecturetimestamp{01:35:52}}
\lecturerresponse{Yes. After expanding about the vacuum, powers of the
quantized field generate vertices. The quartic term contains a four-Higgs
vertex, and the shifted potential also produces cubic interactions.
Knowing their coefficients predicts processes such as
\(hh\to hh\), subject to the usual diagrammatic calculation.}
\end{classroomqa}

\begin{classroomqa}
\audiencequestion{Must a fundamental Lagrangian be polynomial in its
fields, with no power above four?\lecturetimestamp{01:36:09}}
\lecturerresponse{For a perturbatively renormalizable four-dimensional
scalar theory, operators of dimension at most four form the traditional
finite list. Modern effective field theory is less restrictive: higher
dimension operators are allowed and encode short-distance physics, but
their effects are suppressed by powers of the high scale. Thus the
occasional exception is essential to the old rule.}
\end{classroomqa}

\begin{classroomqa}
\audiencequestion{Were self-interaction terms omitted from the discussion
of running couplings?\lecturetimestamp{01:36:44}}
\lecturerresponse{They were included conceptually. Gluon self-interactions
are exactly what reverse the sign relative to QED and produce
antiscreening. Higgs self-interactions likewise contribute to the running
of scalar and other Standard Model parameters.}
\end{classroomqa}

\begin{classroomqa}
\audiencequestion{Are there cross-terms in which different interactions
affect one another's running?\lecturetimestamp{01:37:23}}
\lecturerresponse{Yes. At higher loop order, beta functions contain mixed
products of gauge, Yukawa, and scalar couplings. Perturbation theory remains
useful while those dimensionless couplings are sufficiently small, because
successive loop orders are correspondingly suppressed.}
\end{classroomqa}

\begin{classroomqa}
\audiencequestion{What exactly was meant by saying that \(F\) is finely
tuned?\lecturetimestamp{01:37:44}}
\lecturerresponse{The complete answer is deferred in the lecture. The
issue is the instability of a light scalar scale against much larger
ultraviolet contributions, summarized by
\eqref{eq:l10-higgs-quadratic-sensitivity}. Supersymmetry was one proposed
way to enforce cancellations, but the hierarchy problem itself does not
select a confirmed solution.}
\end{classroomqa}

\begin{classroomqa}
\audiencequestion{Can the renormalization group itself be explained now?
\lecturetimestamp{01:37:56}}
\lecturerresponse{Not in the remaining minutes. Its operative content has
already appeared: change the resolution, let the couplings flow, and
preserve fixed physical predictions. A full treatment belongs to the next
stage of the course.}
\end{classroomqa}

That is where the sequence properly leaves us. The Higgs mechanism has
made the Standard Model internally coherent: it preserves the chiral gauge
symmetry, gives masses to weak bosons and fermions, and predicts physical
Higgs interactions. Renormalization then teaches us how those interactions
change with scale and hints that gauge forces may share a unified origin.
Gravity supplies a second high scale by an independent argument. None of
those successes explains why the Higgs vacuum value is so far below both.
The final lecture therefore does not close the subject; it identifies the
small number on which the next subject must concentrate.
